% 图 53.5 —— 由 `checks/emit_hand.py` 自动拼出（probe 精测 + prims2 图元分解）
%%
% ⚠ 这是**稿子不是成品**：跑 `checks/checkhand.py 53.5` 看指标 + 看叠合图。
%%
%% ⚠ 待核：
%   · 本图 probe 报出阴影族 1 个 —— **本脚本不生成阴影**（要照原书手写）；prims2 的 strokes 里可能含阴影线，核对时留意别画重
%   · 可疑字形 3 项（空心圈 ⊙ 常被认成字母 o）—— 见文件末
%%
\documentclass[12pt]{standalone}
\usepackage{fontspec}
\setsansfont{Arial}
\newfontfamily\mfallback{DejaVu Sans}
\usepackage{tikz}
\usetikzlibrary{arrows.meta}
\begin{document}
\begin{tikzpicture}[x=1pt, y=-1pt, line width=5.0pt, line cap=round, line join=round]

  %% ════ 填充区（prims2 的 fills）════
  \fill (447.0,395.0) -- (445.0,404.0) -- (448.0,407.0) -- (450.0,405.0) -- (450.0,397.0) -- cycle;
  \fill (294.0,489.0) -- (291.0,490.0) -- (288.0,496.0) -- (289.0,500.0) -- (293.0,497.0) -- (295.0,493.0) -- cycle;

  %% ════ 直线段（prims2 的 strokes）════
  \draw[line width=5.0pt] (486.0,331.0) -- (508.0,354.0);
  \draw[line width=5.0pt] (360.0,248.0) -- (361.0,122.0);
  \draw[line width=5.0pt] (283.0,506.0) -- (278.0,513.0);
  \draw[line width=6.0pt] (238.0,470.0) -- (238.0,464.0);
  \draw[line width=5.0pt] (256.0,510.0) -- (277.0,514.0);
  \draw[line width=3.0pt] (150.0,254.0) -- (145.0,257.0);
  \draw[line width=5.0pt] (487.0,504.0) -- (482.0,498.0);
  \draw[line width=3.8pt] (319.2,377.0) -- (316.0,375.0);
  \draw[line width=5.0pt] (280.0,422.0) -- (308.0,424.0);
  \draw[line width=5.0pt] (524.0,400.0) -- (509.0,355.0);
  \fill (502.9,357.0) -- (515.1,353.0) -- (515.1,373.2) -- cycle;
  \draw[line width=5.0pt] (488.0,505.0) -- (507.0,500.0);
  \draw[line width=5.0pt] (527.0,316.0) -- (618.0,317.0);
  \draw[line width=5.0pt] (234.0,413.0) -- (244.0,369.0);
  \draw[line width=4.0pt] (619.0,318.0) -- (615.0,331.0);
  \draw[line width=5.0pt] (458.0,458.0) -- (455.0,451.0);
  \draw[line width=5.0pt] (306.0,387.0) -- (310.0,376.0);
  \draw[line width=5.0pt] (305.0,388.0) -- (279.0,417.0);
  \draw[line width=7.0pt] (483.0,411.0) -- (490.0,407.2);
  \draw[line width=5.0pt] (212.0,501.0) -- (237.0,471.0);
  \draw[line width=3.0pt] (212.0,501.0) -- (209.0,501.0);
  \draw[line width=5.0pt] (212.0,501.0) -- (255.0,510.0);
  \draw[line width=5.0pt] (270.0,428.0) -- (239.0,463.0);
  \draw[line width=5.0pt] (302.0,473.0) -- (298.0,480.0);
  \draw[line width=5.3pt] (250.0,365.0) -- (249.0,360.0);
  \draw[line width=6.0pt] (447.0,397.0) -- (448.0,403.0);
  \draw[line width=6.3pt] (309.0,375.0) -- (303.0,374.0);
  \draw[line width=5.0pt] (448.0,387.0) -- (450.0,377.0);
  \draw[line width=4.0pt] (360.0,316.0) -- (361.0,250.0);
  \draw[line width=5.0pt] (454.0,442.0) -- (452.0,433.0);
  \draw[line width=5.0pt] (341.0,343.0) -- (359.0,318.0);
  \draw[line width=5.3pt] (341.0,343.0) -- (336.0,344.0);
  \draw[line width=5.0pt] (21.0,283.0) -- (42.2,242.8);
  \draw[line width=5.0pt] (507.0,355.0) -- (453.0,369.0);
  \draw[line width=4.0pt] (761.0,317.0) -- (708.0,317.0);
  \draw[line width=4.0pt] (311.0,375.0) -- (315.0,375.0);
  \draw[line width=5.0pt] (361.0,317.0) -- (500.8,317.0);
  \draw[line width=3.0pt] (500.8,317.0) -- (503.0,319.0);
  \draw[line width=5.0pt] (249.0,360.0) -- (265.0,339.0);
  \draw[line width=5.0pt] (136.0,591.0) -- (208.0,502.0);
  \draw[line width=6.0pt] (504.0,319.0) -- (526.0,316.0);
  \draw[line width=5.0pt] (458.0,350.0) -- (462.0,344.0);
  \draw[line width=5.0pt] (238.0,471.0) -- (255.0,509.0);
  \draw[line width=5.0pt] (531.0,244.0) -- (555.0,275.7);
  \draw[line width=5.0pt] (533.0,308.0) -- (527.0,315.0);
  \draw[line width=5.0pt] (304.0,456.0) -- (303.0,462.0);
  \draw[line width=4.0pt] (308.0,438.0) -- (306.0,447.0);
  \draw[line width=5.0pt] (302.0,365.0) -- (303.0,373.0);
  \draw[line width=5.0pt] (361.0,120.0) -- (362.0,43.0);
  \draw[line width=5.0pt] (251.0,366.0) -- (302.0,374.0);
  \draw[line width=5.0pt] (449.0,419.0) -- (482.0,412.0);
  \draw[line width=5.0pt] (316.0,374.0) -- (335.0,345.0);
  \draw[line width=5.0pt] (452.0,368.0) -- (453.0,360.0);
  \draw[line width=5.0pt] (470.0,485.0) -- (476.0,491.0);
  \draw[line width=4.0pt] (503.0,320.0) -- (486.0,330.0);
  \draw[line width=5.7pt] (501.0,414.0) -- (501.0,419.0);
  \draw[line width=5.0pt] (705.0,317.0) -- (620.0,317.0);
  \draw[line width=5.0pt] (334.0,344.0) -- (266.0,338.0);
  \draw[line width=7.1pt] (244.0,369.0) -- (250.0,366.0);
  \draw[line width=7.0pt] (279.0,423.0) -- (272.0,427.0);
  \draw[line width=5.0pt] (675.5,236.5) -- (701.0,282.0);
  \draw[line width=4.0pt] (701.0,282.0) -- (706.0,316.0);
  \draw[line width=5.0pt] (464.0,474.0) -- (461.0,469.0);
  \draw[line width=5.0pt] (238.0,463.0) -- (234.0,415.0);

  %% ════ 曲线（prims2 的 curves —— 三段式，坐标同为 (y,x)）════
  \draw[line width=5.0pt] (486.0,505.0) .. controls (420.6,522.5) and (346.1,525.4) .. (278.0,514.0);
  \draw[line width=5.0pt] (451.0,369.0) .. controls (409.2,378.5) and (364.0,377.1) .. (319.2,377.0);
  \draw[line width=4.0pt] (619.0,316.0) .. controls (620.5,294.7) and (604.0,279.4) .. (590.0,265.0);
  \draw[line width=5.0pt] (502.0,420.0) .. controls (510.8,417.1) and (511.6,425.6) .. (507.2,430.8);
  \draw[line width=6.0pt] (507.2,430.8) .. controls (499.3,435.5) and (498.1,426.4) .. (501.0,421.0);
  \draw[line width=7.0pt] (490.0,407.2) .. controls (491.7,412.3) and (486.8,415.5) .. (483.0,412.0);
  \draw[line width=7.0pt] (278.0,418.0) .. controls (272.4,416.3) and (269.7,421.0) .. (271.0,426.0);
  \draw[line width=5.0pt] (485.0,330.0) .. controls (438.7,338.9) and (390.3,346.9) .. (341.0,343.0);
  \draw[line width=5.0pt] (359.0,249.0) .. controls (300.7,252.4) and (242.8,260.4) .. (194.0,292.0);
  \draw[line width=3.0pt] (624.0,368.0) .. controls (631.8,362.6) and (632.1,347.4) .. (624.0,342.0);
  \draw[line width=4.0pt] (207.0,501.0) .. controls (113.9,477.2) and (-14.0,398.2) .. (21.0,283.0);
  \draw[line width=5.0pt] (42.2,242.8) .. controls (121.9,149.2) and (243.5,124.5) .. (360.0,121.0);
  \draw[line width=5.0pt] (188.0,248.0) .. controls (179.0,258.6) and (167.9,269.6) .. (170.0,284.8);
  \draw[line width=5.0pt] (170.0,284.8) .. controls (189.3,319.2) and (233.4,324.6) .. (265.0,337.0);
  \draw[line width=7.0pt] (249.0,360.0) .. controls (243.0,359.3) and (239.8,364.2) .. (244.0,369.0);
  \draw[line width=5.0pt] (508.0,499.0) .. controls (525.2,470.4) and (528.3,434.6) .. (524.0,401.0);
  \draw[line width=5.0pt] (509.0,500.0) .. controls (594.9,478.1) and (703.2,420.2) .. (707.0,318.0);
  \draw[line width=5.0pt] (555.0,275.7) .. controls (557.2,290.2) and (544.5,301.9) .. (533.0,308.0);
  \draw[line width=4.0pt] (525.0,400.0) .. controls (543.9,394.2) and (563.4,387.9) .. (579.0,374.0);
  \draw[line width=5.0pt] (448.0,420.0) .. controls (403.3,428.3) and (355.8,429.0) .. (310.0,424.0);
  \draw[line width=4.0pt] (116.0,283.0) .. controls (99.1,300.4) and (98.3,328.9) .. (109.4,349.4);
  \draw[line width=4.0pt] (109.4,349.4) .. controls (139.6,389.4) and (187.8,405.0) .. (233.0,414.0);
  \draw[line width=5.0pt] (362.0,249.0) .. controls (424.4,249.2) and (484.9,256.8) .. (537.0,289.0);
  \draw[line width=5.0pt] (362.0,121.0) .. controls (477.1,123.3) and (594.5,147.5) .. (675.5,236.5);

  %% ════ 标签（probe 直接给的 \node 行）════
  %% ⚠ 全部补了 fill=white：文字在最上层，否则与曲线交叉干扰
     \node[anchor=center,inner sep=0pt,fill=white,font=\fontsize{24.5}{30.6}\selectfont] at (363.4,26.3) {\textsf{\textbf{\textit{Z}}}};   % box(354,17,16,17)
     \node[anchor=center,inner sep=0pt,fill=white,font=\fontsize{52.0}{65.0}\selectfont] at (556.9,245.4) {\textsf{\textbf{.}}};   % box(549,242,8,6)
     \node[anchor=center,inner sep=0pt,fill=white,font=\fontsize{47.4}{59.3}\selectfont] at (577.6,255.9) {\textsf{\textbf{.}}};   % box(570,253,8,5)
     \node[anchor=center,inner sep=0pt,fill=white,font=\fontsize{56.1}{70.2}\selectfont] at (134.2,270.0) {\textsf{\textbf{.}}};   % box(126,266,8,7)
     \node[anchor=center,inner sep=0pt,fill=white,font=\fontsize{29.8}{37.3}\selectfont] at (779.7,324.2) {\textsf{\textbf{\textit{y}}}};   % box(771,311,17,23)
     \node[anchor=center,inner sep=0pt,fill=white,font=\fontsize{30.0}{37.5}\selectfont] at (616.8,354.2) {\textsf{\textbf{\textit{b}}}};   % box(607,342,17,22)
     \node[anchor=center,inner sep=0pt,fill=white,font=\fontsize{28.5}{35.6}\selectfont] at (533.0,354.5) {\textsf{\textbf{\textit{f}}}};   % box(527,343,11,22)
     \node[anchor=center,inner sep=0pt,fill=white,font=\fontsize{32.6}{40.7}\selectfont] at (546.0,357.5) {\textsf{\textbf{(}}};   % box(541,343,9,28)
     \node[anchor=center,inner sep=0pt,fill=white,font=\fontsize{34.9}{43.6}\selectfont] at (587.6,357.6) {\textsf{\textbf{×}}};   % box(579,347,16,16)
     \node[anchor=center,inner sep=0pt,fill=white,font=\fontsize{43.8}{54.8}\selectfont] at (295.3,364.5) {\textsf{\textbf{′}}};   % box(292,348,8,11)
     \node[anchor=center,inner sep=0pt,fill=white,font=\fontsize{30.0}{37.6}\selectfont] at (561.5,358.5) {\textsf{\textbf{\textit{a}}}};   % box(553,349,15,16)
     \node[anchor=center,inner sep=0pt,fill=white,font=\fontsize{33.0}{41.2}\selectfont] at (228.6,367.1) {\textsf{\textbf{\textit{a}}}};   % box(219,357,17,17)
     \node[anchor=center,inner sep=0pt,fill=white,font=\fontsize{32.9}{41.2}\selectfont] at (142.3,403.3) {\textsf{\textbf{\textit{C}}}};   % box(131,391,21,24)
     \node[anchor=center,inner sep=0pt,fill=white,font=\fontsize{16.3}{20.3}\selectfont] at (311.4,406.3) {\textsf{\textbf{1}}};   % box(307,400,5,12)
     \node[anchor=center,inner sep=0pt,fill=white,font=\fontsize{24.7}{30.9}\selectfont] at (157.6,417.9) {\textsf{\textbf{2}}};   % box(151,409,12,17)
     \node[anchor=center,inner sep=0pt,fill=white,font=\fontsize{31.5}{39.3}\selectfont] at (323.8,469.8) {\textsf{\textbf{\textit{C}}}};   % box(313,458,20,23)
     \node[anchor=center,inner sep=0pt,fill=white,font=\fontsize{31.1}{38.9}\selectfont] at (616.3,482.5) {\textsf{\textbf{\textit{D}}}};   % box(604,471,21,22)
     \node[anchor=center,inner sep=0pt,fill=white,font=\fontsize{23.0}{28.7}\selectfont] at (340.7,483.9) {\textsf{\textbf{1}}};   % box(334,476,8,15)
     \node[anchor=center,inner sep=0pt,fill=white,font=\fontsize{28.3}{35.4}\selectfont] at (123.9,601.9) {\textsf{\textbf{\textit{x}}}};   % box(115,592,16,16)

  % 钉死画布：否则超出的图元/控制点会把页面撑大，1:1 回比就失准
  \pgfresetboundingbox
  \path[use as bounding box] (0,0) rectangle (797,619);
\end{tikzpicture}
\end{document}

% ⚠ 待核清单（probe 拿不准，人眼过一遍）：
%   · 未识别/跳过：⑤ 跳过/未识别的字形
%   · 未识别/跳过：box(144,253,8,7)
%   · 未识别/跳过：box(624,342,9,28)
